\chapter*{数学符号}
\BookTopBookmark{数学符号}{book-mathematical-notation}
\markboth{数学符号}{}

\begin{fullwidth}
  下表汇总全书反复使用的记号。各章若赋予符号更具体的含义，会在首次使用时另行说明。

  \begingroup
  \small
  \setlength{\tabcolsep}{5pt}
  \renewcommand{\arraystretch}{1.14}
  \vspace{5mm}
  \noindent{\heiti\bfseries 对象与线性代数\par}
  \smallskip
  \noindent\begin{tabularx}{\linewidth}{@{}p{31mm}p{56mm}>{\raggedright\arraybackslash}X@{}}
    \toprule
    对象 & 常用记号 & 含义 \\
    \midrule
    数集
      & $\mathbb{N},\ \mathbb{Z},\ \mathbb{R},\ \mathbb{R}_{+}$
      & 自然数集、整数集、实数集和非负实数集 \\
    集合与关系
      & $x\in A,\ A\subseteq B,\ A\cup B,\ A\cap B,\lvert A\rvert$
      & 元素属于集合、子集、并集、交集和有限集合的基数 \\
    标量
      & $a,\ x,\ \lambda\in\mathbb{R}$
      & 单个数值；小写斜体字母通常表示标量 \\
    向量
      & $\bm{x},\ \bm{w}\in\mathbb{R}^{d}$
      & $d$ 维向量；加粗小写字母表示向量 \\
    矩阵
      & $\bm{A}\in\mathbb{R}^{m\times n},\ \bm{A}_{i,j}$
      & $m$ 行 $n$ 列矩阵及其第 $i$ 行、第 $j$ 列元素 \\
    数据与样本
      & $\mathcal{D}=\{(\bm{x}_i,y_i)\}_{i=1}^{n}$
      & 含 $n$ 个输入标签对的数据集；下标 $i$ 表示样本索引 \\
    切片与迭代
      & $\bm{A}_{i,:},\ \bm{A}_{:,j},\ \bm{w}^{(t)}$
      & 矩阵的第 $i$ 行、第 $j$ 列，以及第 $t$ 次迭代的向量 \\
    Kronecker delta
      & $\delta_{i,j}$
      & 当 $i=j$ 时取 $1$，否则取 $0$ \\
    矩阵变换
      & $\bm{A}^{\top},\ \bm{A}^{-1},\ \bm{A}^{\dagger},\
        \bm{I}_{d},\ \operatorname{tr}(\bm{A})$
      & 转置、逆、伪逆、$d$ 阶单位矩阵和矩阵的迹 \\
    代数运算
      & $\langle\bm{x},\bm{y}\rangle,\ \bm{x}^{\top}\bm{A}\bm{x},\
        \sum_i a_i,\ \prod_i a_i$
      & 内积、二次型、求和与连乘 \\
    函数与模型
      & $f\colon\mathcal{X}\to\mathcal{Y},\ h_{\bm{w}}(\bm{x})$
      & 从输入空间到输出空间的函数，以及参数为 $\bm{w}$ 的模型 \\
    范数与距离
      & $\lVert\bm{x}\rVert_p,\ d(\bm{x},\bm{y})$
      & $p$ 范数与两对象之间的距离；具体定义由上下文给出 \\
    \bottomrule
  \end{tabularx}

  \medskip
  \noindent\begin{minipage}{\linewidth}
  {\heiti\bfseries 微积分与优化\par}
  \smallskip
  \noindent\begin{tabularx}{\linewidth}{@{}p{31mm}p{56mm}>{\raggedright\arraybackslash}X@{}}
    \toprule
    对象 & 常用记号 & 含义 \\
    \midrule
    导数与梯度
      & $\dfrac{\mathrm{d}f}{\mathrm{d}x},\
        \nabla_{\bm{w}}L(\bm{w})$
      & 一元导数，以及标量值函数对参数向量的梯度 \\
    Jacobian 与 Hessian
      & $\nabla_{\bm{x}}\bm{f},\
        \nabla^{2}_{\bm{x}}f$
      & 向量值函数的一阶偏导矩阵，以及标量值函数的二阶偏导矩阵 \\
    渐近阶
      & $f(n)=\mathcal{O}(g(n))$
      & 函数增长率的渐近上界 \\
    最优化
      & $\min_{\bm{w}}L(\bm{w}),\
        \bm{w}^{*}\in\argmin_{\bm{w}}L(\bm{w})$
      & 最小目标值与取得最小值的参数；最大化使用 $\max$ 和 $\argmax$ \\
    \bottomrule
  \end{tabularx}
  \end{minipage}

  \medskip
  \noindent\begin{minipage}{\linewidth}
  {\heiti\bfseries 概率与信息论\par}
  \smallskip
  \noindent\begin{tabularx}{\linewidth}{@{}p{31mm}p{56mm}>{\raggedright\arraybackslash}X@{}}
    \toprule
    对象 & 常用记号 & 含义 \\
    \midrule
    随机变量与实现
      & $X,\ \bm{X},\ x,\ \bm{x},\ p_X(x)$
      & 随机标量、随机向量及其实现值；$p_X$ 表示概率质量或密度函数 \\
    联合与条件概率
      & $p(x,y),\ p(x\mid y)$
      & 随机变量实现值的联合概率与条件概率 \\
    边缘概率
      & $p(x)=\sum_y p(x,y),\
        p(x)=\int p(x,y)\,\mathrm{d}y$
      & 分别对离散变量求和或对连续变量积分得到的边缘概率 \\
    独立性与抽样
      & $X\perp Y,\ X\perp Y\mid Z,\ X_i\overset{\mathrm{i.i.d.}}{\sim}p$
      & 独立、给定 $Z$ 条件独立，以及独立同分布抽样 \\
    矩与相关性
      & $\mathbb{E}[X],\ \operatorname{Var}(X),\
        \operatorname{Cov}(X,Y)$
      & 期望、方差和协方差 \\
    常用分布
      & $\operatorname{Bernoulli}(p),\
        \operatorname{Cat}(\bm{\pi}),\
        \mathcal{N}(\bm{\mu},\bm{\Sigma})$
      & 伯努利分布、类别分布和多元高斯分布 \\
    事件指示函数
      & $\mathbb{I}\{A\}$
      & 事件 $A$ 成立时取 $1$，否则取 $0$ \\
    信息量
      & $H(p),\ H(p,q),\ \operatorname{KL}(p\Vert q)$
      & 熵、交叉熵和分布 $p$ 相对于 $q$ 的KL散度 \\
    概率正比关系
      & $p(\bm{x})\propto q(\bm{x})$
      & 忽略归一化常数时的概率正比关系 \\
    \bottomrule
  \end{tabularx}
  \end{minipage}
  \endgroup
\end{fullwidth}
